% Generated draft: requires mathematical and source review.
\documentclass[11pt]{article}
\usepackage[margin=25mm]{geometry}
\usepackage{fontspec}
\setmainfont{Latin Modern Roman}
\usepackage{amsmath,amssymb,mathtools,unicode-math}
\setmathfont{Latin Modern Math}
\usepackage{enumitem,needspace,xurl,hyperref}
\setlength{\parindent}{0pt}
\setlength{\parskip}{5pt}
\setlength{\emergencystretch}{3em}
\newcommand{\N}{\mathbb{N}}
\newcommand{\Z}{\mathbb{Z}}
\newcommand{\Q}{\mathbb{Q}}
\newcommand{\R}{\mathbb{R}}
\newcommand{\C}{\mathbb{C}}
\newcommand{\F}{\mathbb{F}}
\newcommand{\A}{\mathbb{A}}
\newcommand{\PP}{\mathbb{P}}
\newcommand{\E}{\mathbb{E}}
\newcommand{\eps}{\varepsilon}
\newcommand{\dd}{\,\mathrm{d}}
\begin{document}
\section*{UnsolvedMath ID: 1900015}
\textbf{Source problem code:} catalogue entry
\subsection{Definitions and mathematical statement}
The source literally asks: ``Find frequencies on $n$-dimensional continued fractions with the highest relative frequencies.'' The repeated word ``frequencies'' does not identify the objects to be maximized. The title and supplied research summary interpret the target as \emph{face types of sails having highest relative frequencies}; that is a proposed interpretation, not a silent correction of the displayed source sentence.

Under that interpretation, a mathematical formulation would first specify a family of $n$-dimensional continued fractions, an equivalence relation defining face types $\tau$, and a sampling or invariant-measure convention. For example, if a specified exhaustion gives finite face samples $\mathcal F_R$, the relative frequency of type $\tau$ would be
\[
p(\tau)=\lim_{R\to\infty}
\frac{|\{F\in\mathcal F_R:F\text{ has type }\tau\}|}{|\mathcal F_R|},
\]
provided the denominator is nonzero eventually and the limit exists. The conditional maximization target would then be to determine all types satisfying
\[
p(\tau)=\sup_{\sigma}p(\sigma),
\]
including whether this supremum is attained. A measure-based frequency would instead require a specified probability space and event for each type. Neither construction, nor their equivalence, is supplied by the record; these displays explain the missing data rather than define the source's intended frequency.

The supplied summary reports that the classical $n=2$ setting is substantially understood through the Gauss map and generalizations, with face-type distributions studied, while for $n\geq3$ statistics are sparse and no complete answer about highest frequencies was located. It labels this OPEN-TRIAGE with moderate confidence. The meaning of $n$ as ambient dimension versus continued-fraction dimension is not fixed in the source sentence, and the summary's $n=2$ convention is retained only as its own context.
\subsection{Short English statement}
The likely target is to identify the most frequent sail face types. The supplied sentence is malformed and does not define a frequency measure, sampling process, or face equivalence, so the maximization cannot yet be given a unique rigorous meaning. The summary reports partial classical understanding and no complete higher-dimensional answer.
\subsection{Sources}
\begin{itemize}
\item Local UnsolvedMath registry, permanent ID 1900015, exact source code catalogue entry, ``Geometry of Continued Fractions --- Sail statistics''; status open, classification OPEN-TRIAGE.
\item Karpenkov, \emph{Open problems in geometry of continued fractions} (2017), Problem 15, \url{https://arxiv.org/abs/1712.01450}. The record identifies extraction from public arXiv source TeX, accessed July 29, 2026, and says the source presented the item as open.
\item The supplied moderate-confidence review mentions Arnold and Karpenkov, \emph{Multidimensional continued fractions}, and unspecified ``frequency'' papers for face distributions, without precise bibliographic details or theorem statements. It describes $n=2$ as classical and reports no complete higher-dimensional maximization result; this is a supplied assessment, not an independent verification.
\end{itemize}
\end{document}
